\section{Scalable SPIN: A Growing Outer and One-Bit Route}
\label{sec:structured-distance}

Set $p=1$ in the regional encoder of Section~\ref{sec:regional-spin}.
We call its outer groups \emph{rows}. One sampled constituent is reused
across all rows, and its length grows logarithmically. The route sends
distinct bits of each row to distinct regions. A fixed-width recursive
inner then mixes the routed stream. We will prove a distance guarantee
over the sampled outer, route, and inner setup while retaining linear
encoding work. The proof must account for both sparse and dense sets of
active rows.

\subsection{Scalable Constituent Family}
\label{sec:structured-scalable}

The scalable outer begins with an even base length $a$ and an injective
rate-$1/2$ map $H_a:\F_2^{a/2}\longrightarrow\F_2^a$.
Its weight enumerator
\[
  W_a(z):=\sum_{v\in\F_2^{a/2}}z^{\wt(H_a(v))}
\]
must be known exactly or bounded by a certified envelope. For any multiple
$b$ of $a$, let $H_{a,b}$ be the direct sum of $b/a$ copies of $H_a$.
Setup samples two independent permutations $\tau_1,\tau_2\gets S_b$ and
constructs the rate-$1/2$ BAA constituent
\begin{equation}
  \mathcal O_{a,b}
  :=
  \Acc_b\circ\tau_2\circ
  \Acc_b\circ\tau_1\circ H_{a,b}.
  \label{eq:ba3-outer}
\end{equation}
It samples $\mathcal O_{a,b}$ once and reuses that realization in every outer
position. This is the chosen-block form of the BAA construction
\cite{KolesnikovEtAl2026BA,cryptoeprint:2026/1903}.

The base length $a$ is a construction parameter. A distance certificate must
fix a particular $H_a$, its spectrum, and an admissible schedule for $b$.
Changing $a$ or $H_a$ changes the outer enumerator and may change the certified
schedule. It does not change the route or the structured inner.

\paragraph{The Asymptotic Instantiation.}
For the theorem below, set $a:=24$ and let
$H_{24}:\F_2^{12}\longrightarrow\F_2^{24}$ encode the extended binary Golay code. Fix a positive integer $m$ and set
$L_m:=128m$. Let $b_m$ be the least positive multiple of $24$ satisfying
\begin{equation}
  b_m\ge \frac{39}{4}\log_2(L_mb_m),
  \qquad
  N_m:=L_mb_m,
  \qquad
  K_m:=N_m/2.
  \label{eq:structured-native-schedule}
\end{equation}
Then $b_m=\frac{39}{4}\log_2N_m+O(1)$, and consecutive native lengths have relative
gap $o(1)$.

Use $\mathcal O_{24,b_m}$ from \eqref{eq:ba3-outer} with local dimension
$k_m=b_m/2$. Setup reuses this sampled
map in all $L_m$ diagonal positions:
$G_{m,1}=\cdots=G_{m,L_m}:=\mathcal O_{24,b_m}$.
The expectation over $(\tau_1,\tau_2)$ remains outside powers of the
constituent enumerator; this matters when counting several active rows.

Setup samples the bit-level group and region shuffles of
Section~\ref{sec:regional-spin}, independently of $(\tau_1,\tau_2)$.
Fix $(t,s)=(64,22)$ and the $s=22$ map
$A:\F_2^{22}\to\F_2^{64}$ specified in Appendix~\ref{app:shared-inner-map}.
Set $C=A^\top$. The map $A$ is injective, and every column of $C$ is
nonzero. Its exact image spectrum and the spectrum of $\ker C$ enter
the distance calculation. For each step, independently sample
$M_i\gets\mathrm{GL}(22,2)$, independently of the outer and route.
For every fixed nonzero state $q$, $M_iq$ is uniform over the nonzero
states. All these choices belong to one setup fixed for every message.
The same $s=22$ inner map and update law appear in the larger finite
Turbo profiles of Section~\ref{sec:turbo-spin}.

\begin{theorem}[Scalable SPIN]
\label{thm:structured-spin-scalable}
For the setup above,
\begin{equation}
  \Prob\left[
    d_{\min}(E_m)\le\lfloor0.11N_m\rfloor
  \right]
  =o(1)
  \qquad(m\to\infty).
  \label{eq:structured-spin-distance}
\end{equation}
The probability is over the two BAA permutations, all route permutations,
and the independent uniform state maps. Every realized code has rate $1/2$. Its ordinary
and transposed encoders use $O(N_m)$ bit operations.
\end{theorem}

The argument has three parts. First, control the realized spectrum of the
one sampled outer. Next, bound the route--inner failure probability for
each number of active blocks. Finally, sum those bounds before applying
the first-moment theorem. Appendix~\ref{app:structured-scalable-certificate}
gives the spectrum calculations, transfer entries, and numerical witnesses
used below.
The block constant $39/4$ is the value certified here; a smaller certified
constant would change the schedule, not the construction.

\subsection{Selecting the Shared Outer}

For $p=1$, the outer proof interface is its local weight enumerator.
For a realized local map $G_j$, define
\begin{equation}
  P_{G_j}(z)
  :=
  \sum_{v\in\F_2^k\setminus\{0\}}z^{\wt(G_j(v))}.
  \label{eq:local-constituent-enumerator}
\end{equation}
The variable $z$ records weight: the coefficient of $z^w$ counts the
nonzero block messages whose encoded output has weight $w$.
These coefficients are ordinary integer counts, not elements of $\F_2$.
We write $[z^h]F(z)$ for the coefficient of $z^h$ in a polynomial $F(z)$.
If one deterministic map $G$ is repeated, the number of messages with exactly
$Q$ active blocks and outer weight $h$ is
\begin{equation}
  \binom LQ [z^h]P_G(z)^Q.
  \label{eq:repeated-constituent-enumerator}
\end{equation}
Here $\binom LQ$ chooses the $Q$ nonzero message-block positions.
Expanding $P_G(z)^Q$ chooses one nonzero message for each of these blocks;
multiplication adds their output weights in the exponent of $z$.
The coefficient $[z^h]$ therefore counts exactly the choices with total
outer weight $h$. For example, if $P_G(z)=2z^2+z^4$, then
$[z^6]P_G(z)^2=4$: either block can supply weight two, with two choices
for its message, while the other supplies weight four.

For independently sampled, identically distributed local maps, expectation
replaces $P_G$ in \eqref{eq:repeated-constituent-enumerator} by
$\E[P_{G_1}]$. If one sampled map is reused in all positions, the correct
quantity is instead
\begin{equation}
  \binom LQ\E\bigl[[z^h]P_G(z)^Q\bigr].
  \label{eq:shared-constituent-enumerator}
\end{equation}
In particular, an expectation over one sampled constituent must not be moved
inside the $Q$th power.

The distance proof need not know how the local encoder is implemented. It may
consume an exact enumerator, an expected enumerator under the declared setup,
or a coefficientwise certified envelope. Minimum distance alone is generally
insufficient because nearby low-weight shells can have much larger
multiplicity.

Suppress the family index $m$, and write $\mathcal O_b:=\mathcal O_{24,b}$.
Let $A_{\mathcal O_b}(w)$ count its nonzero words of weight $w$, and let
$\overline A_b(w):=\E_{\tau_1,\tau_2}A_{\mathcal O_b}(w)$.
The exact formula in Lemma~\ref{lem:structured-exact-ba-spectrum} composes
the Golay enumerator with two accumulator weight transitions.
It supplies both the tail bound and the spectrum envelope.
Let $\mathcal W:=[0.104,0.896]$. Lemma~\ref{lem:structured-ba-tails}
proves
\begin{equation}
  \sum_{w/b\notin\mathcal W}\overline A_b(w)=o(1).
  \label{eq:structured-ba-tail}
\end{equation}

\begin{lemma}[One-sample selection]
\label{lem:structured-one-sample-selection}
There is an event $\mathcal G_b$, determined by the two BAA permutations, such
that
\begin{equation}
  \Prob[\mathcal G_b]=1-o(1).
  \label{eq:ba-good-event}
\end{equation}
On $\mathcal G_b$, every nonzero constituent word has relative weight in
$\mathcal W$, and
\begin{equation}
  A_{\mathcal O_b}(w)
  \le b^2\overline A_b(w)
  \qquad(0\le w\le b).
  \label{eq:ba-selected-spectrum}
\end{equation}
\end{lemma}

\begin{proof}
The expected number of nonzero tail words equals the left side of
\eqref{eq:structured-ba-tail}. Markov's inequality shows that no such word
exists with probability $1-o(1)$. For each remaining weight $w$, another
application of Markov's inequality gives
\[
  \Prob[A_{\mathcal O_b}(w)>b^2\overline A_b(w)]\le b^{-2}.
\]
A union bound over at most $b+1$ weights costs $O(1/b)$. Intersecting the two
events proves the lemma.
\end{proof}

The $o(1)$ loss from Lemma~\ref{lem:structured-one-sample-selection} is
incurred once.
After setup samples $\mathcal O_b$, the construction reuses that realization
at every outer position.

Fix any realization satisfying $\mathcal G_b$. All subsequent bounds are
uniform over this choice; the remaining randomness consists of the
independent route permutations and uniform state maps.
The spectrum calculation also gives a concave function
$\widehat a_{\rm BA}$ on $\mathcal W$ such that
\[
 A_{\mathcal O_b}(w)
 \le \exp\{b\widehat a_{\rm BA}(w/b)+O(\ln b)\}.
\]
Appendix~\ref{app:structured-scalable-certificate} derives this majorant
from the exact expected spectrum, including the uniform remainder.
Selection lets us multiply bounds on the realized spectrum across rows.
It does not replace the expectation of a power by a power of the mean.

For $0\le d\le N$ and $1\le Q\le L$, let $Z_{d,Q}$ count nonzero messages
with exactly $Q$ active outer blocks and output weight at most $d$.
Our goal is to show that
$\sum_{Q=1}^L\E[Z_{\lfloor0.11N\rfloor,Q}\mid\mathcal G_b]=o(1)$.

\subsection{From Row Weights to an iid Comparison}

View an outer word as an $L$-tuple of rows in $\F_2^b$. For the framework in
Section~\ref{sec:framework}, take its class to be the ordered row-weight
profile $c(u):=(\wt(u_1),\ldots,\wt(u_L))$.
This finite class records the information preserved by the independent block
shuffles. The analysis below constructs a uniform route--inner failure
envelope for every such profile. It then aggregates profiles by occupation and
mean row weight where those coarser parameters suffice.

Fix $Q$ active outer positions. Let their realized output weights be
$w_1,\ldots,w_Q$, and set
\[
  x_j:=\frac{w_j}{b},
  \qquad
  \alpha:=\frac QL,
  \qquad
  x:=\frac1Q\sum_{j=1}^Qx_j,
  \qquad
  q:=\alpha x.
\]
When $\alpha$ has a positive lower bound, the next lemma reduces the routing
law to a comparison that depends only on $q$, and the concavity step
\eqref{eq:structured-ba-jensen} reduces the outer spectrum to $(\alpha,x)$.
We use binary relative entropy with natural logarithms,
\[
 D_{\rm KL}(u\Vert v)
 :=u\ln(u/v)+(1-u)\ln((1-u)/(1-v)),
\]
with the usual continuous endpoint conventions.

\begin{lemma}[Route domination]
\label{lem:structured-route-domination}
Let $F$ be any nonnegative function of the routed word, viewed as $b$
regions of $L$ bits.
Conditioned on the row weights $w_1,\ldots,w_Q$, Lemma~\ref{lem:packet-route-symmetry}
reduces the bit-level route to uniform supports. The route then
satisfies
\begin{equation}
  \E_{\rm route}[F]
  \le
  (b+1)^Q(L+1)^b
  \E_{\operatorname{Ber}(q)^{Lb}}[F].
  \label{eq:structured-route-domination}
\end{equation}
If $\alpha$ is bounded below by a positive constant, the multiplicative
factor in \eqref{eq:structured-route-domination} is $\exp(o(N))$.
\end{lemma}

\begin{proof}
For row $j$, begin with $b$ independent Bernoulli-$x_j$ variables and
condition their sum to equal $w_j$. The value $w_j$ is a mode of the
binomial distribution. Hence the inverse conditioning probability is at
most $b+1$. Applying this reduction to all active rows costs at most
$(b+1)^Q$.

Before the region shuffle, one fixed region contains independent Bernoulli
variables whose average success probability is $q$. Let $S$ be their sum.
For every $0\le k\le L$, the Chernoff bound for a Poisson-binomial sum gives
\[
  \Prob[S=k]
  \le
  \exp\{-L D_{\rm KL}(k/L\Vert q)\}.
\]
The method-of-types lower bound gives
\[
  \Prob[\operatorname{Bin}(L,q)=k]
  \ge
  \frac1{L+1}
  \exp\{-L D_{\rm KL}(k/L\Vert q)\}.
\]
After a uniform region shuffle, both laws are uniform on every weight-$k$
slice. The routed region law is therefore pointwise at most $L+1$ times the
iid Bernoulli-$q$ law. Multiplying over the $b$ independent region shuffles
proves \eqref{eq:structured-route-domination}.

Finally,
\[
  \ln((b+1)^Q(L+1)^b)
  =Q\ln(b+1)+b\ln(L+1)=o(Lb)
\]
when $Q/L$ is bounded below by a positive constant and $b,L\to\infty$.
\end{proof}

Concavity of $\widehat a_{\rm BA}$ now reduces the selected outer spectrum
to the mean row weight:
\begin{equation}
  \frac1Q\sum_{j=1}^Q\widehat a_{\rm BA}(x_j)
  \le
  \widehat a_{\rm BA}(x).
  \label{eq:structured-ba-jensen}
\end{equation}
The region shuffles are essential for Lemma~\ref{lem:structured-route-domination}.
The deterministic transpose alone spreads each row, but it does not erase
the row's coordinate pattern.

\subsection{Weighted Inner Bound and Positive Occupation}

We first bound the inner for iid Bernoulli-$\beta$ input bits and a tilt
$0<z<1$. Put $D:=2^{22}-1$ and
\[
 g_0:=1-\beta+\beta z,\qquad g_1:=\beta+(1-\beta)z,
 \qquad f(q):=g_0^{64-\wt(Aq)}g_1^{\wt(Aq)}\quad(q\ne0).
\]
Write $F:=\max_{q\ne0}f(q)$ and
$U:=D^{-1}\sum_{q\ne0}f(q)$. Let
\[
 H:=\sum_{x:Cx=0}\Pr_\beta[X=x]z^{\wt(x)},
 \qquad u:=\Pr_\beta[CX\ne0],
\]
where $X$ is one Bernoulli-$\beta$ step. The image spectrum of $A$
and the kernel spectrum of $C$ determine these quantities.
A weighted state measure is bounded by three coordinates: mass at zero,
arbitrary live mass, and uniform live mass. The independent uniform
mixer bounds their one-step evolution by
\begin{equation}
 T(\beta,z):=
 \begin{pmatrix}
 H&g_0^{64}-H&0\\
 \min(F,u)/D&0&F\\
 \min(U,u)/D&0&U
 \end{pmatrix}.
 \label{eq:shared-inner-transfer}
\end{equation}
For a live state $q$, the mixer sends $q$ uniformly to a nonzero state.
A prescribed successor therefore receives weighted mass at most
$f(q)/D$. Reaching zero additionally requires $CX\ne0$. These facts
give the last two rows of \eqref{eq:shared-inner-transfer}; the first
row separates inputs in $\ker C$. Induction gives
\[
 \E_\beta[z^{\wt(Y)}]\le e_0^\top T(\beta,z)^{N/64}\1_3.
\]
Appendix~\ref{app:shared-inner-proof} proves the envelope for the
full state process; it does not assume uniform feedback syndromes.

For $\alpha:=Q/L\ge10^{-4}$ and mean active-row weight $x\in\mathcal W$,
Lemma~\ref{lem:structured-route-domination} compares the routed input
with iid Bernoulli-$\alpha x$ input. Changing that reference density to
$py$, with $p,y\in(0,1)$, costs at most
$\exp\{N[D_{\rm KL}(\alpha\Vert p)+\alpha D_{\rm KL}(x\Vert y)]\}$.
If $T(py,e^{-\lambda})\boldsymbol v\le\rho\boldsymbol v$ for a
positive vector $\boldsymbol v$, the bad-message exponent per output
bit is at most
\begin{equation}
 \alpha\widehat a_{\rm BA}(x)
 +D_{\rm KL}(\alpha\Vert p)+\alpha D_{\rm KL}(x\Vert y)
 +\frac{\ln\rho}{64}+0.11\lambda+o(1).
 \label{eq:shared-dense-exponent}
\end{equation}
For fixed $p,y,\lambda$ and one affine outer-majorant segment, the
expression without $o(1)$ is convex in $(\alpha,\alpha x)$.
The certificate partitions $[10^{-4},1]\times\mathcal W$ into rational
boxes and checks their four transformed vertices with outward arithmetic.
Its largest upper exponent is below $-1.00147\cdot10^{-7}$.
The finite witness set has bounded moment prefactors. Thus this
occupation range contributes at most
$L\exp\{-10^{-7}N+o(N)\}=o(1)$.

\subsection{Growing Sparse Occupation}

The factor $(L+1)^b$ in Lemma~\ref{lem:structured-route-domination} is
too large when $Q/L$ tends to zero. Compare each active row with a fair
$b$-bit row instead. The outer spectrum bounds the cost of this
comparison by $\exp\{b((\ln2)/2+\ell+\epsilon_b)\}$ per row, where
$\ell=0.01281$ and $\epsilon_b=O(\ln b/b)$.
Generate Bernoulli marks at density $p=(8/5)\alpha$ and condition each
region to have $Q$ marks. Before conditioning, the routed input bits
are iid Bernoulli-$\beta$ with $\beta=(4/5)\alpha$. The inverse
conditioning probability is at most
$8\sqrt Q\exp\{L D_{\rm KL}(\alpha\Vert p)\}$ per region.

Set $z:=1-(8/5)\alpha$ and
$\boldsymbol v(\alpha):=(1,(1+24\alpha)/1024,1/1024)^\top$.
The exact rational spectrum check gives
\begin{equation}
 T(\beta,z)\boldsymbol v(\alpha)
 \le(1-48\alpha)\boldsymbol v(\alpha)
 \qquad(0<\alpha\le10^{-4}).
 \label{eq:shared-sparse-collatz}
\end{equation}
The moment prefactor is at most $1024$. Combining this contraction,
outer-word counts, mark conditioning, and the schedule
\eqref{eq:structured-native-schedule} yields
\begin{equation}
 \E[Z_{\lfloor0.11N\rfloor,Q}\mid\mathcal G_b]
 \le e^{-0.006Qb}
 \qquad(4096\le Q\le10^{-4}L)
 \label{eq:shared-sparse-final}
\end{equation}
The cutoff makes the remainder uniform even when $Q$ grows slowly.
Appendix~\ref{app:shared-inner-proof} gives the exact contraction and
all conditioning remainders.

\subsection{Fixed Occupation}

Fix $Q<4096$ while $L$ grows. A region then contains at most $Q$
active bits. Its shuffle places them at distinct uniform slots.
With probability $1-O_Q(1/L)$, their steps are separated by empty
steps and avoid the region boundaries. From a live state, one empty
step makes the successor exactly uniform nonzero. The output weight
of that first step still depends on the entering state, so the proof
compares the \emph{full-state} regional kernels before multiplying
across regions.

Let $p_0:=2^{21}/D$ and $z_L:=e^{-\theta/L}$. For an empty gap of
$g\ge1$ steps, the live-to-live weighted kernel from $q$ to $q'$ is
\[
 \frac{z_L^{\wt(Aq)}}{D}
 \left(D^{-1}\sum_{v\ne0}z_L^{\wt(Av)}\right)^{g-1}.
\]
An isolated input bit activates a zero state because every column of
$C$ is nonzero. From a uniform live state it returns to zero with
probability $1/D$. Thus the limiting zero/live impulse matrix is
$P_0=\left(\begin{smallmatrix}0&1\\1/D&1-1/D\end{smallmatrix}\right)$.
For $a$ bits placed in a region, let $F_{a,L}(z_L)$ be its full-state
weighted kernel and $K_a(\theta)$ the limiting two-state kernel.
Appendix~\ref{app:shared-inner-proof}
defines $K_a$ by a simplex integral and proves the componentwise
full-state comparison
\[
 F_{a,L}(z_L)\le(1+C_{Q,\theta}/L)J K_a(\theta)R
 \qquad(0\le a\le Q).
\]
Here $J$ classifies the entering state and $R$ distributes an exiting
live class uniformly over nonzero states. The comparison holds for
every boundary state. Across $b=O(\log L)$ regions its total cost is
$1+o(1)$, including when coefficient extraction fixes row weights.

For $Q=1,2$, the resulting first-moment exponents per active-row bit
are below $-0.109827$ and $-0.109812$. For fixed $3\le Q<4096$,
retaining each row weight and applying a weighted matrix norm gives
\[
 \E[Z_{\lfloor0.11N\rfloor,Q}\mid\mathcal G_b]
 \le\exp\{-0.0008685Qb+o_Q(b)\}.
\]
The finite set of fixed-$Q$ terms therefore tends to zero.

\subsection{Completing the Argument}

\begin{proof}[Proof of Theorem~\ref{thm:structured-spin-scalable}]
Condition on an arbitrary shared outer realization in $\mathcal G_b$.
All route and inner bounds above hold uniformly over that choice.
The fixed-occupation terms vanish in a finite sum.
Equation~\eqref{eq:shared-sparse-final} gives a geometric sum for
$4096\le Q\le10^{-4}L$. The positive-occupation contribution is
$L\exp\{-10^{-7}N+o(N)\}$. Hence
$\sum_{Q=1}^L\E[Z_{\lfloor0.11N\rfloor,Q}\mid\mathcal G_b]=o(1)$.
Markov's inequality and $\Prob[\mathcal G_b^c]=o(1)$ prove the
distance claim.

The outer is injective, and Lemma~\ref{lem:structured-inner-bijective}
proves inner bijectivity for every setup. The Golay base circuits,
two accumulators, outer permutations, and route use $O(N)$ work.
Each fixed-width state update has constant bit-operation cost, so the
inner and its transpose also use $O(N)$ work.
\end{proof}
